\section{Electrical Test System}
\label{sec:model}
\subsection{Feeder and Load Representation}
The generated case, \code{SEIAN_LV_Timed_Switching}, comprises buses G01 and
N02--N06, seven three-phase switchable series R/L branches, and five balanced
Y-connected resistive loads. G01 is supplied by a Thevenin source configured
at 0.4 kV line-to-line RMS and 50 Hz, with a source resistance setting of
$0.02\unit{\Omega}$. The declared source setting and implemented branch values
are illustrative; they are not a short-circuit calibration against a measured
transformer or utility supply.

Five branches form the initial radial tree; N02--N05 and N04--N06 are normally
open ties. Opening N02--N03, N03--N04, and N03--N05 isolates N03. Closing both
ties then supplies the healthy downstream buses through
G01--N02--N05--N06--N04, as shown in Figure~\ref{fig:feeder}.

\begin{figure}[tb]
\centering\includegraphics[width=\textwidth]{feeder.pdf}
\caption{Original schematic of the modeled electrical connectivity before and
after restoration. Lines represent three-phase branches, not communication
links or geographical cable routing. Node positions are diagrammatic.}
\label{fig:feeder}
\end{figure}

\begin{table}[tb]
\centering\small
\caption{Implemented per-phase branch parameters, extracted from the saved
PSCAD build report. These values are illustrative, not measured cable data.}
\label{tab:lines}
\begin{tabular}{@{}lrrl@{}}
\toprule
Branch & $R_e$ ($\Omega$) & $L_e$ ($\mu$H) & Initial state \\
\midrule
G01--N02 & 0.05000 & 47.75 & Closed \\
N02--N03 & 0.07500 & 71.62 & Closed \\
N03--N04 & 0.08571 & 81.85 & Closed \\
N05--N06 & 0.10000 & 95.49 & Closed \\
N03--N05 & 0.10000 & 95.49 & Closed \\
N02--N05 & 0.12000 & 114.59 & Open (tie) \\
N04--N06 & 0.12000 & 114.59 & Open (tie)
\\
\bottomrule
\end{tabular}
\end{table}

The builder assigns resistance from a nominal branch-capacity field. For a
numerical capacity entry $q_e$ in kW, its heuristic is
\begin{equation}
 R_e = \frac{3}{q_e}\unit{\Omega},\qquad
 L_e = \frac{0.3R_e}{2\pi f},\qquad f=50\unit{Hz}.
\end{equation}
The constant 3 is an implementation sizing heuristic with the implied unit
conversion, not a conductor law. Consequently, \code{capacity_kw} should
not be interpreted as an enforced thermal rating. Table~\ref{tab:lines}
reports the actual generated impedances directly.

The loads require similarly careful interpretation. The source builder
contains a scalar named \code{LOAD_PF}, set to $0.98$, but uses it to multiply
the target active power when calculating resistance. It does not introduce a
reactive load or enforce a 0.98 electrical power factor. If $P_i^{\mathrm{set}}$
is expressed in watts and $V_{LL,0}=400\unit{V}$, the implemented phase
resistance is
\begin{equation}
 R_{i,\phi}=
 \frac{(V_{LL,0}/\sqrt{3})^2}{0.98 P_i^{\mathrm{set}}/3}
 =\frac{V_{LL,0}^{2}}{0.98 P_i^{\mathrm{set}}}.
\label{eq:load}
\end{equation}
Thus the nominal-voltage resistive demand is $0.98P_i^{\mathrm{set}}$ and
varies approximately with the square of applied balanced RMS voltage. These
are not constant-power loads. Table~\ref{tab:loads} records both the original
sizing labels and the resulting nominal-voltage demands.

\begin{table}[tb]
\centering\small
\caption{Resistive load sizing. The demand column refers to 400-V operation,
not the measured demand at the loaded bus voltage.}
\label{tab:loads}
\begin{tabular}{@{}lrrr@{}}
\toprule
Bus & Sizing label (kW) & Demand at 400 V (kW) & $R_{i,\phi}$ ($\Omega$) \\
\midrule
N02 & 8 & 7.84 & 20.40816 \\
N03 & 8 & 7.84 & 20.40816 \\
N04 & 10 & 9.80 & 16.32653 \\
N05 & 8 & 7.84 & 20.40816 \\
N06 & 6 & 5.88 & 27.21088
\\
\bottomrule
\end{tabular}
\end{table}

\subsection{Fault and Numerical Configuration}
The physical-fault scenario uses a native three-phase-to-ground fault at N03,
with an on-resistance of $0.05\unit{\Omega}$ and an off-resistance of
$10^6\unit{\Omega}$. Its programmed inception is 4.8 s, duration 0.4 s,
and clearing time 5.2 s. All incident N03 breakers are commanded open at
5.0 s and both ties close at 6.0 s. Fault activation uses the native
\code{tpflt} and \code{tfaultn} components; it is not synthesized by manually
changing a plotted voltage trace. Finite off-resistance produces small leakage
when the fault is inactive. The native resistive-switch abstraction does not
model fault-arc dynamics \cite{pscadFaults}.

PSCAD advances the EMT solution with a $25\unit{\mu s}$ time step and records
research signals every $50\unit{\mu s}$. A duration-$T$ experiment therefore
contains $T/(50\unit{\mu s})+1$ recorded samples per channel when both endpoints
are included. Startup is included in the saved trajectories but excluded from
the interruption metric as described below. Time-step convergence, frequency
dependence of line parameters, grounding calibration, unbalanced demand,
inverter dynamics, and transformer detail have not been established for this
case.
